\section{Conclusion}
\label{sec:conclusion}

This paper formulated Windows malware model development as auditable evidence
allocation under partial observability. The decision stage combines early static
evidence with event-indexed dynamic evidence; the adaptation stage assigns a
limited labeling budget; and the audit stage tests the attribution ranking used
in model inspection and acquisition. The formulation provides one scientific
flow while preserving component-specific estimands and datasets.

On the 1,922-sample MalBehavD-V1 development cohort, Adaptive P reduced FPR from
12.45\% to 11.16\% relative to fixed weighting and reduced Brier by 0.00306,
with paired intervals excluding zero. Full PID obtained AUC comparable to
logistic stacking and a lower FPR at its independently selected operating point,
with a small TPR trade-off. Integral and derivative terms did not add consistent
value over P within 1--20 event prefixes. The evidence therefore favors a
simpler proportional adaptation and treats richer PID variants as hypotheses for
external timestamped evaluation.

On the identity-audited 2,000-sample EMBER subset, $U+H$ produced the strongest
ROC/AP learning-curve areas. $U+V$ reached the predeclared target with the fewest
labels, while \xfull used 492 labels, 92 fewer than Random. Gains over
Uncertainty were smaller and metric dependent. Across 25 fitted models,
top-$|\mathrm{SHAP}|$ deletion produced much larger AUC decreases than equal-
budget random deletion under mean and median replacement, while moderate
stability revealed dependence on data partitions.

These results support component-wise adaptive fusion, explanation-guided sample
prioritization, and quantitative attribution auditing for generic Windows
malware. They do not establish ransomware detection, causal PE explanations,
external generalization, or an end-to-end integrated system. The next evidence
gate is a frozen external or temporal paired cohort with common operating targets
and real timestamps, followed by stronger AL diversity baselines, conditional
deletion controls, and a factorial acquisition--fusion experiment on one cohort.
